\documentclass[12pt]{article}
\usepackage[utf8]{inputenc}
\usepackage[T1]{fontenc}
\usepackage{hyperref}
\title{Quechan Language Resources for K-12 Curriculum Development}
\author{Compiled by Claude Sonnet 5.5 (claude-sonnet-5-5) for Andr\'as Kornai}
\date{}
\begin{document}
\maketitle

\section{Introduction}
Quechan (Kwats\'aan; ISO 639-3 yum; Glottocode quec1382) is a Yuman language of the Fort Yuma Quechan Tribe in California and Arizona. It is severely endangered. Sources give widely differing speaker counts, for example about 60 in 2020 on Wikipedia and 290 in a census figure on Omniglot.

\section{Primary Educational Resources}
\begin{itemize}
\item \emph{Kwats\'aan Iiy\'aa Mattkuu'\'eeyk! Learn the Quechan Language!}, a book developed for tribal language classes, covering the most frequently taught words for beginner and intermediate learners. The UBC RelLex entry is at \url{https://wiki.ubc.ca/Documentation:RelLex/Kwats%C3%A1an_Iiy%C3%A1a_Mattkuu%27%C3%A9eyk!_Learn_the_Quechan_Language!}. A PDF appears to be hosted by the tribe at \url{https://quechantribe.com/documents/download/47}.
\end{itemize}

\section{Linguistic and Reference Materials}
\begin{itemize}
\item \emph{Kwats\'aan liv\'aa: A Multi-Generational Dictionary of the Quechan Language}, which builds on A.M. Halpern's work and tribal members' work with Amy Miller: \url{https://wiki.ubc.ca/Documentation:RelLex/Kwats%C3%A1an_liv%C3%A1a:_A_Multi-Generational_Dictionary_of_the_Quechan_Language}
\item UBC RelLex page on Quechan: \url{https://wiki.ubc.ca/Documentation:RelLex/Quechan}
\item Open Book Publishers PDF with an introduction by Amy Miller describing the forthcoming Quechan dictionary and Halpern's story collection: \url{https://books.openbookpublishers.com/10.11647/obp.0037.pdf}
\item Wikipedia overview: \url{https://en.wikipedia.org/wiki/Quechan_language}
\item Native-languages.org Quechan page, listing lessons and linguistic descriptions: \url{https://www.native-languages.org/quechan.htm}
\item Omniglot page on Quechan, including orthography: \url{https://www.omniglot.com/writing/quechan.htm}
\end{itemize}

\section{Religious and Community Resources}
The Quechan Tribe website (\url{https://quechantribe.com/documents/download/47}) is the community source found. No scripture translations were found.

\section{Digital and Audio Resources}
No audio, video or software resources were identified in the searches.

\section{Curriculum Development Context}
There is an existing classroom-oriented book and a community-based dictionary, which is a better base than for most endangered languages. With so few speakers, any curriculum depends on tribal language programme leadership and elder involvement. A practical orthography was developed in the 1990s with elders.

\end{document}